\subsection*{Cauchy Sequences}

\begin{exposition}
Cauchy sequences measure internal stabilization. Instead of comparing the
terms to an already named limit, the Cauchy condition compares sufficiently
late terms to each other. Over the real numbers this internal criterion is
equivalent to convergence, and that equivalence is one of the central
sequence forms of completeness.
\end{exposition}

% ---------------------------------------------------------
% Definition --- Cauchy Sequence
% ---------------------------------------------------------

\begin{definitionbox}{Definition (Cauchy Sequence)}
\begin{definition}[Cauchy Sequence]
\label{def:cauchy-sequence}
Let $(x_n)$ be a real sequence. The sequence $(x_n)$ is a
\textbf{Cauchy sequence} if
\[
  \forall \varepsilon>0 \;\exists N\in\mathbb{N}\;\forall m,n\ge N
  \;\bigl(|x_m-x_n|<\varepsilon\bigr).
\]
\end{definition}
\end{definitionbox}

\begin{remark*}[Standard quantified statement]
\[
\begin{aligned}
&(x_n)\text{ is Cauchy}\\
&\quad\Longleftrightarrow\quad
\forall \varepsilon>0 \;\exists N\in\mathbb{N}\;\forall m,n\ge N
\;\bigl(|x_m-x_n|<\varepsilon\bigr).
\end{aligned}
\]
\end{remark*}

\begin{remark*}[Definition predicate reading]
\[
\operatorname{CauchySequence}(x_n)
\coloneqq
\forall \varepsilon>0 \;\exists N\in\mathbb{N}\;\forall m,n\ge N
\;\bigl(|x_m-x_n|<\varepsilon\bigr).
\]
\end{remark*}

\begin{remark*}[Negated quantified statement]
\[
\begin{aligned}
&(x_n)\text{ is not Cauchy}\\
&\quad\Longleftrightarrow\quad
\exists \varepsilon>0 \;\forall N\in\mathbb{N}\;\exists m,n\ge N
\;\bigl(|x_m-x_n|\ge\varepsilon\bigr).
\end{aligned}
\]
\end{remark*}

\begin{remark*}[Negation predicate reading]
\[
\neg\operatorname{CauchySequence}(x_n)
\Longleftrightarrow
\exists \varepsilon>0 \;\forall N\in\mathbb{N}\;\exists m,n\ge N
\;\bigl(|x_m-x_n|\ge\varepsilon\bigr).
\]
\end{remark*}

\begin{remark*}[Interpretation]
The Cauchy condition says that the tails of the sequence eventually have
arbitrarily small internal spread. It does not name a candidate limit; it
only asserts that late terms become mutually close.
\end{remark*}

\begin{dependencies}
\begin{itemize}
  \item \hyperref[def:real-sequence]{Real Sequence}
\end{itemize}
\end{dependencies}

% ---------------------------------------------------------
% Theorem --- Convergent Sequences Are Cauchy
% ---------------------------------------------------------

\begin{theorembox}{Theorem (Convergent Sequences Are Cauchy)}
\begin{theorem}[Convergent Sequences Are Cauchy]
\label{thm:convergent-sequences-are-cauchy}
Let $(x_n)$ be a real sequence. If $(x_n)$ converges, then $(x_n)$ is
Cauchy.

\smallskip
\noindent
\hyperref[prf:convergent-sequences-are-cauchy]{\textit{Go to proof.}}
\end{theorem}
\end{theorembox}

\begin{remark*}[Standard quantified statement]
\[
\begin{aligned}
&\forall (x_n)\subseteq\mathbb{R}\\
&\quad
\Bigl(
  \bigl(\exists L\in\mathbb{R}\;\forall \varepsilon>0\;\exists N\in\mathbb{N}
  \;\forall n\ge N\;( |x_n-L|<\varepsilon )\bigr)\\
&\qquad\Longrightarrow
  \forall \varepsilon>0\;\exists N\in\mathbb{N}\;\forall m,n\ge N
  \;( |x_m-x_n|<\varepsilon )
\Bigr).
\end{aligned}
\]
\end{remark*}

\begin{remark*}[Definition predicate reading]
\[
\forall (x_n)\subseteq\mathbb{R}\;
\Bigl(
  \exists L\in\mathbb{R}\;\operatorname{ConvergentSequence}(x_n,L)
  \Longrightarrow
  \operatorname{CauchySequence}(x_n)
\Bigr).
\]
\end{remark*}

\begin{remark*}[Interpretation]
Once every sufficiently late term is close to the same limit, any two
sufficiently late terms are close to each other. This theorem is the easy
half of the Cauchy criterion.
\end{remark*}

\begin{dependencies}
\begin{itemize}
  \item \hyperref[def:convergent-sequence]{Convergence of a Sequence}
  \item \hyperref[def:cauchy-sequence]{Cauchy Sequence}
\end{itemize}
\end{dependencies}

% ---------------------------------------------------------
% Theorem --- Cauchy Sequences Are Bounded
% ---------------------------------------------------------

\begin{theorembox}{Theorem (Cauchy Sequences Are Bounded)}
\begin{theorem}[Cauchy Sequences Are Bounded]
\label{thm:cauchy-sequences-are-bounded}
Every Cauchy real sequence is bounded.

\smallskip
\noindent
\hyperref[prf:cauchy-sequences-are-bounded]{\textit{Go to proof.}}
\end{theorem}
\end{theorembox}

\begin{remark*}[Standard quantified statement]
\[
\begin{aligned}
&\forall (x_n)\subseteq\mathbb{R}\\
&\quad
\Bigl(
  \forall \varepsilon>0\;\exists N\in\mathbb{N}\;\forall m,n\ge N
  \;( |x_m-x_n|<\varepsilon )\\
&\qquad\Longrightarrow
  \exists M>0\;\forall n\in\mathbb{N}\;( |x_n|\le M )
\Bigr).
\end{aligned}
\]
\end{remark*}

\begin{remark*}[Definition predicate reading]
\[
\forall (x_n)\subseteq\mathbb{R}\;
\bigl(
  \operatorname{CauchySequence}(x_n)
  \Longrightarrow
  \operatorname{BoundedSequence}(x_n)
\bigr).
\]
\end{remark*}

\begin{remark*}[Interpretation]
A Cauchy sequence has one tightly controlled tail. The finitely many terms
before that tail can be absorbed into a single larger bound.
\end{remark*}

\begin{dependencies}
\begin{itemize}
  \item \hyperref[def:cauchy-sequence]{Cauchy Sequence}
  \item \hyperref[def:bounded-sequence]{Bounded Sequence}
\end{itemize}
\end{dependencies}

% ---------------------------------------------------------
% Theorem --- Cauchy Sequence with a Convergent Subsequence
% ---------------------------------------------------------

\begin{theorembox}{Theorem (Cauchy Sequence with a Convergent Subsequence)}
\begin{theorem}[Cauchy Sequence with a Convergent Subsequence]
\label{thm:cauchy-sequence-with-convergent-subsequence}
Let $(x_n)$ be a Cauchy real sequence. If a subsequence of $(x_n)$
converges to $L\in\mathbb{R}$, then $(x_n)$ converges to $L$.

\smallskip
\noindent
\hyperref[prf:cauchy-sequence-with-convergent-subsequence]{\textit{Go to proof.}}
\end{theorem}
\end{theorembox}

\begin{remark*}[Standard quantified statement]
\[
\begin{aligned}
&\forall (x_n)\subseteq\mathbb{R}\;\forall L\in\mathbb{R}\\
&\quad
\Bigl(
  \bigl[
    \forall \varepsilon>0\;\exists N\in\mathbb{N}\;\forall m,n\ge N
    \;( |x_m-x_n|<\varepsilon )
  \bigr]\\
&\qquad\land
  \bigl[
    \exists \sigma:\mathbb{N}\to\mathbb{N}\;
    \bigl(
      \forall k,\ell\in\mathbb{N}\;(k<\ell\Longrightarrow \sigma(k)<\sigma(\ell))
      \land
      \forall \varepsilon>0\;\exists K\in\mathbb{N}\;\forall k\ge K
      \;( |x_{\sigma(k)}-L|<\varepsilon )
    \bigr)
  \bigr]\\
&\qquad\Longrightarrow
  \forall \varepsilon>0\;\exists N\in\mathbb{N}\;\forall n\ge N
  \;( |x_n-L|<\varepsilon )
\Bigr).
\end{aligned}
\]
\end{remark*}

\begin{remark*}[Definition predicate reading]
\[
\begin{aligned}
&\forall (x_n)\subseteq\mathbb{R}\;\forall L\in\mathbb{R}\\
&\quad
\Bigl(
  \operatorname{CauchySequence}(x_n)
  \land
  \operatorname{SubsequentialLimit}(x_n,L)
  \Longrightarrow
  \operatorname{ConvergentSequence}(x_n,L)
\Bigr).
\end{aligned}
\]
\end{remark*}

\begin{remark*}[Interpretation]
The Cauchy condition prevents the rest of the tail from wandering away from
any convergent subsequence. A single subsequential limit therefore determines
the limit of the whole sequence.
\end{remark*}

\begin{dependencies}
\begin{itemize}
  \item \hyperref[def:cauchy-sequence]{Cauchy Sequence}
  \item \hyperref[def:subsequence]{Subsequence}
  \item \hyperref[def:subsequential-limit]{Subsequential Limit}
  \item \hyperref[def:convergent-sequence]{Convergence of a Sequence}
\end{itemize}
\end{dependencies}

% ---------------------------------------------------------
% Theorem --- Cauchy Criterion for Real Sequences
% ---------------------------------------------------------

\begin{theorembox}{Theorem (Cauchy Criterion for Real Sequences)}
\begin{theorem}[Cauchy Criterion for Real Sequences]
\label{thm:cauchy-criterion-real-sequences}
Let $(x_n)$ be a real sequence. Then $(x_n)$ converges if and only if
$(x_n)$ is Cauchy.

\smallskip
\noindent
\hyperref[prf:cauchy-criterion-real-sequences]{\textit{Go to proof.}}
\end{theorem}
\end{theorembox}

\begin{remark*}[Standard quantified statement]
\[
\begin{aligned}
&\forall (x_n)\subseteq\mathbb{R}\\
&\quad
\Bigl(
  \exists L\in\mathbb{R}\;\forall \varepsilon>0\;\exists N\in\mathbb{N}
  \;\forall n\ge N\;( |x_n-L|<\varepsilon )\\
&\qquad\Longleftrightarrow
  \forall \varepsilon>0\;\exists N\in\mathbb{N}\;\forall m,n\ge N
  \;( |x_m-x_n|<\varepsilon )
\Bigr).
\end{aligned}
\]
\end{remark*}

\begin{remark*}[Definition predicate reading]
\[
\forall (x_n)\subseteq\mathbb{R}\;
\Bigl(
  \exists L\in\mathbb{R}\;\operatorname{ConvergentSequence}(x_n,L)
  \Longleftrightarrow
  \operatorname{CauchySequence}(x_n)
\Bigr).
\]
\end{remark*}

\begin{remark*}[Interpretation]
This theorem is the sequence form of completeness of $\mathbb{R}$. In an
incomplete ambient space, a sequence may be internally stable without
converging inside that space.
\end{remark*}

\begin{dependencies}
\begin{itemize}
  \item \hyperref[thm:convergent-sequences-are-cauchy]{Convergent Sequences Are Cauchy}
  \item \hyperref[thm:cauchy-sequences-are-bounded]{Cauchy Sequences Are Bounded}
  \item \hyperref[thm:bolzano-weierstrass-sequences]{Bolzano-Weierstrass Theorem for Sequences}
  \item \hyperref[thm:cauchy-sequence-with-convergent-subsequence]{Cauchy Sequence with a Convergent Subsequence}
\end{itemize}
\end{dependencies}

% ---------------------------------------------------------
% Theorem --- Cauchy Criterion via Tails
% ---------------------------------------------------------

\begin{theorembox}{Theorem (Cauchy Criterion via Tails)}
\begin{theorem}[Cauchy Criterion via Tails]
\label{thm:cauchy-criterion-via-tails}
Let $(x_n)$ be a real sequence. Then $(x_n)$ is Cauchy if and only if
for every $\varepsilon>0$ there exists $N\in\mathbb{N}$ such that the
$N$-tail has pairwise term distances less than $\varepsilon$.

\smallskip
\noindent
\hyperref[prf:cauchy-criterion-via-tails]{\textit{Go to proof.}}
\end{theorem}
\end{theorembox}

\begin{remark*}[Standard quantified statement]
\[
\begin{aligned}
&(x_n)\text{ is Cauchy}\\
&\quad\Longleftrightarrow\quad
\forall \varepsilon>0\;\exists N\in\mathbb{N}\;\forall p,q\in\mathbb{N}
\;\bigl(|x_{N+p}-x_{N+q}|<\varepsilon\bigr).
\end{aligned}
\]
\end{remark*}

\begin{remark*}[Definition predicate reading]
\[
\operatorname{CauchySequence}(x_n)
\Longleftrightarrow
\forall \varepsilon>0\;\exists N\in\mathbb{N}\;\forall p,q\in\mathbb{N}
\;\bigl(|x_{N+p}-x_{N+q}|<\varepsilon\bigr).
\]
\end{remark*}

\begin{remark*}[Interpretation]
The Cauchy condition is a tail property. Reindexing a sufficiently late tail
does not change the substance of the condition; it only changes the names of
the indices.
\end{remark*}

\begin{dependencies}
\begin{itemize}
  \item \hyperref[def:cauchy-sequence]{Cauchy Sequence}
  \item \hyperref[def:m-tail]{$M$-Tail of a Sequence}
\end{itemize}
\end{dependencies}

% ---------------------------------------------------------
% Theorem --- Cauchy Tail Diameter Criterion
% ---------------------------------------------------------

\begin{theorembox}{Theorem (Cauchy Tail Diameter Criterion)}
\begin{theorem}[Cauchy Tail Diameter Criterion]
\label{thm:cauchy-tail-diameter-criterion}
Let $(x_n)$ be a real sequence. Then $(x_n)$ is Cauchy if and only if
every sufficiently late tail has arbitrarily small pairwise diameter.

\smallskip
\noindent
\hyperref[prf:cauchy-tail-diameter-criterion]{\textit{Go to proof.}}
\end{theorem}
\end{theorembox}

\begin{remark*}[Standard quantified statement]
\[
\begin{aligned}
&(x_n)\text{ is Cauchy}\\
&\quad\Longleftrightarrow\quad
\forall \varepsilon>0\;\exists N_0\in\mathbb{N}\;\forall N\ge N_0\\
&\qquad\qquad
\forall m,n\ge N\;\bigl(|x_m-x_n|<\varepsilon\bigr).
\end{aligned}
\]
\end{remark*}

\begin{remark*}[Definition predicate reading]
\[
\operatorname{CauchySequence}(x_n)
\Longleftrightarrow
\forall \varepsilon>0\;\exists N_0\in\mathbb{N}\;\forall N\ge N_0\;\forall m,n\ge N
\;\bigl(|x_m-x_n|<\varepsilon\bigr).
\]
\end{remark*}

\begin{remark*}[Interpretation]
The tail-diameter form says that once one tail is small, every later tail is
small as well. It records the monotone nature of tail control without
introducing additional notation for diameters.
\end{remark*}

\begin{dependencies}
\begin{itemize}
  \item \hyperref[def:cauchy-sequence]{Cauchy Sequence}
  \item \hyperref[def:m-tail]{$M$-Tail of a Sequence}
\end{itemize}
\end{dependencies}

% ---------------------------------------------------------
% Theorem --- Cauchy Sequences Have Vanishing Successive Differences
% ---------------------------------------------------------

\begin{theorembox}{Theorem (Cauchy Sequences Have Vanishing Successive Differences)}
\begin{theorem}[Cauchy Sequences Have Vanishing Successive Differences]
\label{thm:cauchy-successive-differences-vanish}
If $(x_n)$ is a Cauchy real sequence, then
\[
  |x_{n+1}-x_n|\to 0.
\]

\smallskip
\noindent
\hyperref[prf:cauchy-successive-differences-vanish]{\textit{Go to proof.}}
\end{theorem}
\end{theorembox}

\begin{remark*}[Standard quantified statement]
\[
\begin{aligned}
&\forall (x_n)\subseteq\mathbb{R}\\
&\quad
\Bigl(
  \forall \varepsilon>0\;\exists N\in\mathbb{N}\;\forall m,n\ge N
  \;( |x_m-x_n|<\varepsilon )\\
&\qquad\Longrightarrow
  \forall \varepsilon>0\;\exists N\in\mathbb{N}\;\forall n\ge N
  \;( |x_{n+1}-x_n|<\varepsilon )
\Bigr).
\end{aligned}
\]
\end{remark*}

\begin{remark*}[Definition predicate reading]
\[
\operatorname{CauchySequence}(x_n)
\Longrightarrow
\operatorname{NullSequence}\bigl(|x_{n+1}-x_n|\bigr).
\]
\end{remark*}

\begin{remark*}[Interpretation]
Cauchy control applies to every pair of late indices, so it applies in
particular to consecutive late indices. The converse is false in general:
small successive steps do not prevent accumulated drift.
\end{remark*}

\begin{dependencies}
\begin{itemize}
  \item \hyperref[def:cauchy-sequence]{Cauchy Sequence}
  \item \hyperref[def:null-sequence]{Null Sequence}
\end{itemize}
\end{dependencies}

% ---------------------------------------------------------
% Theorem --- Scalar Multiples of Cauchy Sequences
% ---------------------------------------------------------

\begin{theorembox}{Theorem (Scalar Multiples of Cauchy Sequences)}
\begin{theorem}[Scalar Multiples of Cauchy Sequences]
\label{thm:scalar-multiple-cauchy-sequence}
Let $(x_n)$ be a Cauchy real sequence, and let $\alpha\in\mathbb{R}$.
Then $(\alpha x_n)$ is Cauchy.

\smallskip
\noindent
\hyperref[prf:scalar-multiple-cauchy-sequence]{\textit{Go to proof.}}
\end{theorem}
\end{theorembox}

\begin{remark*}[Standard quantified statement]
\[
\begin{aligned}
&\forall (x_n)\subseteq\mathbb{R}\;\forall \alpha\in\mathbb{R}\\
&\quad
\Bigl(
  (x_n)\text{ is Cauchy}
  \Longrightarrow
  (\alpha x_n)\text{ is Cauchy}
\Bigr).
\end{aligned}
\]
\end{remark*}

\begin{remark*}[Definition predicate reading]
\[
\operatorname{CauchySequence}(x_n)
\Longrightarrow
\operatorname{CauchySequence}(\alpha x_n).
\]
\end{remark*}

\begin{remark*}[Interpretation]
Multiplying all terms by one fixed real number rescales all pairwise tail
distances by the same fixed factor.
\end{remark*}

\begin{dependencies}
\begin{itemize}
  \item \hyperref[def:cauchy-sequence]{Cauchy Sequence}
  \item \hyperref[def:scalar-multiple-sequence]{Scalar Multiple of a Sequence}
\end{itemize}
\end{dependencies}

% ---------------------------------------------------------
% Theorem --- Sums of Cauchy Sequences
% ---------------------------------------------------------

\begin{theorembox}{Theorem (Sums of Cauchy Sequences)}
\begin{theorem}[Sums of Cauchy Sequences]
\label{thm:sum-cauchy-sequences}
Let $(x_n)$ and $(y_n)$ be Cauchy real sequences. Then $(x_n+y_n)$ is
Cauchy.

\smallskip
\noindent
\hyperref[prf:sum-cauchy-sequences]{\textit{Go to proof.}}
\end{theorem}
\end{theorembox}

\begin{remark*}[Standard quantified statement]
\[
\begin{aligned}
&\forall (x_n)\subseteq\mathbb{R}\;\forall (y_n)\subseteq\mathbb{R}\\
&\quad
\Bigl(
  (x_n)\text{ is Cauchy}
  \land
  (y_n)\text{ is Cauchy}
  \Longrightarrow
  (x_n+y_n)\text{ is Cauchy}
\Bigr).
\end{aligned}
\]
\end{remark*}

\begin{remark*}[Definition predicate reading]
\[
\operatorname{CauchySequence}(x_n)
\land
\operatorname{CauchySequence}(y_n)
\Longrightarrow
\operatorname{CauchySequence}(x_n+y_n).
\]
\end{remark*}

\begin{remark*}[Interpretation]
The triangle inequality splits the tail distance of the sum into the two tail
distances already controlled by the original sequences.
\end{remark*}

\begin{dependencies}
\begin{itemize}
  \item \hyperref[def:cauchy-sequence]{Cauchy Sequence}
  \item \hyperref[def:pointwise-sum-sequence]{Pointwise Sum of Sequences}
\end{itemize}
\end{dependencies}

% ---------------------------------------------------------
% Theorem --- Differences of Cauchy Sequences
% ---------------------------------------------------------

\begin{theorembox}{Theorem (Differences of Cauchy Sequences)}
\begin{theorem}[Differences of Cauchy Sequences]
\label{thm:difference-cauchy-sequences}
Let $(x_n)$ and $(y_n)$ be Cauchy real sequences. Then $(x_n-y_n)$ is
Cauchy.

\smallskip
\noindent
\hyperref[prf:difference-cauchy-sequences]{\textit{Go to proof.}}
\end{theorem}
\end{theorembox}

\begin{remark*}[Standard quantified statement]
\[
\begin{aligned}
&\forall (x_n)\subseteq\mathbb{R}\;\forall (y_n)\subseteq\mathbb{R}\\
&\quad
\Bigl(
  (x_n)\text{ is Cauchy}
  \land
  (y_n)\text{ is Cauchy}
  \Longrightarrow
  (x_n-y_n)\text{ is Cauchy}
\Bigr).
\end{aligned}
\]
\end{remark*}

\begin{remark*}[Definition predicate reading]
\[
\operatorname{CauchySequence}(x_n)
\land
\operatorname{CauchySequence}(y_n)
\Longrightarrow
\operatorname{CauchySequence}(x_n-y_n).
\]
\end{remark*}

\begin{remark*}[Interpretation]
The difference theorem is the sum theorem applied after negating one
sequence. It is often the form used to compare two approximating sequences.
\end{remark*}

\begin{dependencies}
\begin{itemize}
  \item \hyperref[thm:sum-cauchy-sequences]{Sums of Cauchy Sequences}
  \item \hyperref[thm:scalar-multiple-cauchy-sequence]{Scalar Multiples of Cauchy Sequences}
  \item \hyperref[def:pointwise-difference-sequence]{Pointwise Difference of Sequences}
\end{itemize}
\end{dependencies}

% ---------------------------------------------------------
% Theorem --- Linear Combinations of Cauchy Sequences
% ---------------------------------------------------------

\begin{theorembox}{Theorem (Linear Combinations of Cauchy Sequences)}
\begin{theorem}[Linear Combinations of Cauchy Sequences]
\label{thm:linear-combination-cauchy-sequences}
Let $(x_n)$ and $(y_n)$ be Cauchy real sequences, and let
$\alpha,\beta\in\mathbb{R}$. Then $(\alpha x_n+\beta y_n)$ is Cauchy.

\smallskip
\noindent
\hyperref[prf:linear-combination-cauchy-sequences]{\textit{Go to proof.}}
\end{theorem}
\end{theorembox}

\begin{remark*}[Standard quantified statement]
\[
\begin{aligned}
&\forall (x_n)\subseteq\mathbb{R}\;\forall (y_n)\subseteq\mathbb{R}\;
 \forall \alpha\in\mathbb{R}\;\forall \beta\in\mathbb{R}\\
&\quad
\Bigl(
  (x_n)\text{ is Cauchy}
  \land
  (y_n)\text{ is Cauchy}
  \Longrightarrow
  (\alpha x_n+\beta y_n)\text{ is Cauchy}
\Bigr).
\end{aligned}
\]
\end{remark*}

\begin{remark*}[Definition predicate reading]
\[
\operatorname{CauchySequence}(x_n)
\land
\operatorname{CauchySequence}(y_n)
\Longrightarrow
\operatorname{CauchySequence}(\alpha x_n+\beta y_n).
\]
\end{remark*}

\begin{remark*}[Interpretation]
The Cauchy real sequences are closed under finite real linear combinations.
This is the additive vector-space part of Cauchy algebra.
\end{remark*}

\begin{dependencies}
\begin{itemize}
  \item \hyperref[def:linear-combination-of-real-sequences]{Linear Combination of Real Sequences}
  \item \hyperref[thm:scalar-multiple-cauchy-sequence]{Scalar Multiples of Cauchy Sequences}
  \item \hyperref[thm:sum-cauchy-sequences]{Sums of Cauchy Sequences}
\end{itemize}
\end{dependencies}

% ---------------------------------------------------------
% Theorem --- Products of Cauchy Sequences
% ---------------------------------------------------------

\begin{theorembox}{Theorem (Products of Cauchy Sequences)}
\begin{theorem}[Products of Cauchy Sequences]
\label{thm:product-cauchy-sequences}
Let $(x_n)$ and $(y_n)$ be Cauchy real sequences. Then $(x_ny_n)$ is
Cauchy.

\smallskip
\noindent
\hyperref[prf:product-cauchy-sequences]{\textit{Go to proof.}}
\end{theorem}
\end{theorembox}

\begin{remark*}[Standard quantified statement]
\[
\begin{aligned}
&\forall (x_n)\subseteq\mathbb{R}\;\forall (y_n)\subseteq\mathbb{R}\\
&\quad
\Bigl(
  (x_n)\text{ is Cauchy}
  \land
  (y_n)\text{ is Cauchy}
  \Longrightarrow
  (x_ny_n)\text{ is Cauchy}
\Bigr).
\end{aligned}
\]
\end{remark*}

\begin{remark*}[Definition predicate reading]
\[
\operatorname{CauchySequence}(x_n)
\land
\operatorname{CauchySequence}(y_n)
\Longrightarrow
\operatorname{CauchySequence}(x_ny_n).
\]
\end{remark*}

\begin{remark*}[Interpretation]
Products require boundedness in addition to tail closeness. The boundedness
is supplied automatically because every Cauchy real sequence is bounded.
\end{remark*}

\begin{dependencies}
\begin{itemize}
  \item \hyperref[thm:cauchy-sequences-are-bounded]{Cauchy Sequences Are Bounded}
  \item \hyperref[def:pointwise-product-sequence]{Pointwise Product of Sequences}
\end{itemize}
\end{dependencies}

% ---------------------------------------------------------
% Theorem --- Reciprocals of Cauchy Sequences
% ---------------------------------------------------------

\begin{theorembox}{Theorem (Reciprocals of Cauchy Sequences)}
\begin{theorem}[Reciprocals of Cauchy Sequences]
\label{thm:reciprocal-cauchy-sequence}
Let $(x_n)$ be a Cauchy real sequence. If there exist $c>0$ and
$N_0\in\mathbb{N}$ such that $|x_n|\ge c$ for every $n\ge N_0$, then
$(1/x_n)$ is Cauchy.

\smallskip
\noindent
\hyperref[prf:reciprocal-cauchy-sequence]{\textit{Go to proof.}}
\end{theorem}
\end{theorembox}

\begin{remark*}[Standard quantified statement]
\[
\begin{aligned}
&\forall (x_n)\subseteq\mathbb{R}\\
&\quad
\Bigl(
  (x_n)\text{ is Cauchy}
  \land
  \exists c>0\;\exists N_0\in\mathbb{N}\;\forall n\ge N_0\;( |x_n|\ge c )\\
&\qquad\Longrightarrow
  (1/x_n)\text{ is Cauchy}
\Bigr).
\end{aligned}
\]
\end{remark*}

\begin{remark*}[Definition predicate reading]
\[
\operatorname{CauchySequence}(x_n)
\land
\exists c>0\;\exists N_0\in\mathbb{N}\;\forall n\ge N_0\;( |x_n|\ge c )
\Longrightarrow
\operatorname{CauchySequence}(1/x_n).
\]
\end{remark*}

\begin{remark*}[Interpretation]
The lower bound away from zero prevents division from amplifying small tail
differences without control. Without it, reciprocals need not preserve Cauchy
behavior.
\end{remark*}

\begin{dependencies}
\begin{itemize}
  \item \hyperref[def:cauchy-sequence]{Cauchy Sequence}
  \item \hyperref[def:reciprocal-sequence]{Reciprocal Sequence}
  \item \hyperref[thm:reciprocal-of-product-on-r]{Reciprocal of a Product on $\mathbb{R}$}
  \item \hyperref[thm:absolute-value-of-reciprocal-r]{Absolute Value of a Reciprocal on $\mathbb{R}$}
\end{itemize}
\end{dependencies}

% ---------------------------------------------------------
% Theorem --- Quotients of Cauchy Sequences
% ---------------------------------------------------------

\begin{theorembox}{Theorem (Quotients of Cauchy Sequences)}
\begin{theorem}[Quotients of Cauchy Sequences]
\label{thm:quotient-cauchy-sequences}
Let $(x_n)$ and $(y_n)$ be Cauchy real sequences. If there exist $c>0$
and $N_0\in\mathbb{N}$ such that $|y_n|\ge c$ for every $n\ge N_0$,
then $(x_n/y_n)$ is Cauchy.

\smallskip
\noindent
\hyperref[prf:quotient-cauchy-sequences]{\textit{Go to proof.}}
\end{theorem}
\end{theorembox}

\begin{remark*}[Standard quantified statement]
\[
\begin{aligned}
&\forall (x_n)\subseteq\mathbb{R}\;\forall (y_n)\subseteq\mathbb{R}\\
&\quad
\Bigl(
  (x_n)\text{ is Cauchy}
  \land
  (y_n)\text{ is Cauchy}
  \land
  \exists c>0\;\exists N_0\in\mathbb{N}\;\forall n\ge N_0\;( |y_n|\ge c )\\
&\qquad\Longrightarrow
  (x_n/y_n)\text{ is Cauchy}
\Bigr).
\end{aligned}
\]
\end{remark*}

\begin{remark*}[Definition predicate reading]
\[
\operatorname{CauchySequence}(x_n)
\land
\operatorname{CauchySequence}(y_n)
\land
\exists c>0\;\exists N_0\in\mathbb{N}\;\forall n\ge N_0\;( |y_n|\ge c )
\Longrightarrow
\operatorname{CauchySequence}(x_n/y_n).
\]
\end{remark*}

\begin{remark*}[Interpretation]
The quotient theorem is the product theorem combined with the reciprocal
theorem for a denominator eventually bounded away from zero.
\end{remark*}

\begin{dependencies}
\begin{itemize}
  \item \hyperref[thm:product-cauchy-sequences]{Products of Cauchy Sequences}
  \item \hyperref[thm:reciprocal-cauchy-sequence]{Reciprocals of Cauchy Sequences}
  \item \hyperref[def:pointwise-quotient-sequence]{Pointwise Quotient of Sequences}
\end{itemize}
\end{dependencies}

% ---------------------------------------------------------
% Theorem --- Absolute Values of Cauchy Sequences
% ---------------------------------------------------------

\begin{theorembox}{Theorem (Absolute Values of Cauchy Sequences)}
\begin{theorem}[Absolute Values of Cauchy Sequences]
\label{thm:absolute-value-cauchy-sequence}
Let $(x_n)$ be a Cauchy real sequence. Then $(|x_n|)$ is Cauchy.

\smallskip
\noindent
\hyperref[prf:absolute-value-cauchy-sequence]{\textit{Go to proof.}}
\end{theorem}
\end{theorembox}

\begin{remark*}[Standard quantified statement]
\[
\forall (x_n)\subseteq\mathbb{R}\;
\Bigl(
  (x_n)\text{ is Cauchy}
  \Longrightarrow
  (|x_n|)\text{ is Cauchy}
\Bigr).
\]
\end{remark*}

\begin{remark*}[Definition predicate reading]
\[
\operatorname{CauchySequence}(x_n)
\Longrightarrow
\operatorname{CauchySequence}(|x_n|).
\]
\end{remark*}

\begin{remark*}[Interpretation]
The reverse triangle inequality bounds the distance between absolute values
by the original distance between terms.
\end{remark*}

\begin{dependencies}
\begin{itemize}
  \item \hyperref[def:cauchy-sequence]{Cauchy Sequence}
  \item \hyperref[def:absolute-value-sequence]{Absolute Value of a Sequence}
\end{itemize}
\end{dependencies}
